\documentclass[11pt]{article}
\usepackage[a4paper,margin=1in]{geometry}
\usepackage{amsmath,amssymb,amsthm}
\usepackage{graphicx}
\usepackage[hidelinks]{hyperref}
\usepackage{microtype}

\newtheorem{theorem}{Theorem}
\newtheorem{corollary}[theorem]{Corollary}
\newcommand{\N}{\mathbb N}
\newcommand{\F}{\mathbb F}
\newcommand{\cB}{\mathcal B}
\newcommand{\cG}{\mathcal G}

\title{A Finite-Dimensional Obstruction Program\\
for Schauder Reorderings of Quasi-Greedy Markushevich Bases}
\author{Automated mathematics run \texttt{fa\_banach\_001}, lane 11}
\date{9 August 2026}

\begin{document}
\maketitle

\begin{center}
\fbox{\parbox{0.91\textwidth}{\textbf{Status.} Candidate partial result,
likely valid, subject to human review.  This packet proves a finite-dimensional
reduction; it does not solve the open Banach-space problem.}}
\end{center}

\section{Source question}

Miguel Berasategui asks in the introduction of
\cite{Berasategui2025} (arXiv:2510.06398v1, PDF page 2):
\begin{quote}
``It is an open problem whether every quasi-greedy Markushevich basis has a
Schauder reordering.''
\end{quote}
The paper works throughout with Banach spaces.  The following image is a crop
rendered from the original arXiv PDF.

\begin{center}
\includegraphics[width=\textwidth]{figures/open_problem_crop.png}
\end{center}

The more recent quasi-Banach counterexample of
Albiac--Ansorena--Berasategui \cite{AAB2025} explicitly says that this question
remains open under the Banach-space restriction.

\section{Definitions}

Let $E$ be a $d$-dimensional Banach space with a normalized basis
$\cB=(e_i)_{i=1}^d$ and biorthogonal functionals $(e_i^*)_{i=1}^d$.
For $A\subset\{1,\dots,d\}$ write
\[
  P_Ax=\sum_{i\in A}e_i^*(x)e_i.
\]
A finite set $A$ is \emph{greedy for $x$} if
\[
  |e_i^*(x)|\ge |e_k^*(x)|
  \qquad(i\in A,\ k\notin A).
\]
The basis is $K$-quasi-greedy if $\|P_Ax\|\le K\|x\|$ for every
$x\in E$ and every greedy set $A$.

Define the \emph{best ordering constant}
\[
  \beta(\cB)
  :=\min_{\sigma\in S_d}\ \max_{1\le m\le d}
       \left\|P_{\{\sigma(1),\ldots,\sigma(m)\}}\right\|.
\]
Thus $\beta(\cB)$ measures the smallest basis constant obtainable by
reordering the finite basis.

\section{Proof intuition}

Greedy selection is local with respect to an $\ell_2$-sum: a global greedy
set meets every finite block in a greedy set.  Consequently uniform
quasi-greediness of the blocks passes to their union without any loss in the
constant.  Reordering behaves in the opposite direction.  Every ordering of
the union induces an ordering on each block, and a global prefix realizes
each induced block prefix after restriction to that block.  Therefore a block
whose best ordering constant is large forces a large global partial-sum
projection in every ordering.  A sequence of such blocks gives the desired
infinite obstruction.

\section{Finite-dimensional block reduction}

\begin{theorem}[Block obstruction]
Fix $K,M\ge1$.  For each $j\in\N$, let $E_j$ be a finite-dimensional
Banach space with a normalized basis
$\cB_j=(e_{j,r})_{r=1}^{d_j}$ and biorthogonal functionals
$(e_{j,r}^*)_{r=1}^{d_j}$.  Assume
\begin{enumerate}
\item $\sup_{j,r}\|e_{j,r}^*\|\le M$;
\item every $\cB_j$ is $K$-quasi-greedy; and
\item $\beta(\cB_j)\to\infty$.
\end{enumerate}
Then, in the Banach space
\[
  X=\left(\bigoplus_{j=1}^{\infty}E_j\right)_{\ell_2},
\]
the union $\cB=\bigcup_j\cB_j$, in any enumeration, is a normalized,
$M$-bounded, $K$-quasi-greedy Markushevich basis.  No permutation of $\cB$
is a Schauder basis of $X$.
\end{theorem}

\begin{proof}
Embed each $e_{j,r}$ in the $j$th coordinate of $X$, and similarly embed
$e_{j,r}^*$ in $X^*=(\bigoplus_jE_j^*)_{\ell_2}$.  The resulting systems are
biorthogonal.  Finite block-supported vectors are dense in $X$, and if all
embedded coordinate functionals vanish at $x=(x_j)_j$, then every $x_j=0$.
Thus the union is a Markushevich basis.  It is normalized and $M$-bounded by
hypothesis.

Let $x=(x_j)_j\in X$, and let $A$ be a finite greedy set for $x$ relative to
the union.  Put $A_j=A\cap\{(j,r):1\le r\le d_j\}$.  If $A_j$ is nonempty,
then for every selected $r$ and every unselected $s$ in the same block, the
global greedy inequality gives
\[
  |e_{j,r}^*(x_j)|\ge |e_{j,s}^*(x_j)|.
\]
Hence $A_j$ is a greedy set for $x_j$ in $\cB_j$; the empty case contributes
zero.  Therefore
\[
 \|P_Ax\|_X^2
   =\sum_j\|P_{A_j}x_j\|_{E_j}^2
   \le K^2\sum_j\|x_j\|_{E_j}^2
   =K^2\|x\|_X^2.
\]
This proves that the union is $K$-quasi-greedy.

Now fix an arbitrary permutation $\pi$ of the union.  Its restriction to the
$j$th block induces a permutation $\sigma_j\in S_{d_j}$.  By the definition
of $\beta(\cB_j)$, there exists $m_j\in\{1,\ldots,d_j\}$ such that
\[
 \left\|P_{\{\sigma_j(1),\ldots,\sigma_j(m_j)\}}\right\|
 \ge \beta(\cB_j).
\]
Let $N_j$ be the position in the global order $\pi$ of the coordinate
$e_{j,\sigma_j(m_j)}$.  At time $N_j$, precisely the first $m_j$ coordinates
of the induced order $\sigma_j$ have appeared from block $j$.  Consequently,
the global initial projection $S_{N_j}^{\pi}$, restricted to the isometric
copy of $E_j$, is exactly the displayed block projection.  Hence
\[
  \|S_{N_j}^{\pi}\|\ge\beta(\cB_j).
\]
Because $\beta(\cB_j)\to\infty$, the initial projections associated with
$\pi$ are not uniformly bounded.  Thus the union is not a Schauder basis in
the order $\pi$.  Since $\pi$ was arbitrary, no Schauder reordering exists.
\end{proof}

\begin{corollary}[Necessary finite-dimensional consequence of a positive answer]
If every quasi-greedy Markushevich basis in a Banach space admits a Schauder
reordering, then for every $K,M\ge1$ there is $C(K,M)<\infty$ such that
\[
  \beta(\cB)\le C(K,M)
\]
for every normalized finite-dimensional $K$-quasi-greedy basis whose
coordinate functionals have norm at most $M$.
\end{corollary}

\begin{proof}
If no such $C(K,M)$ existed, choose one finite basis for each integer $j$
with $\beta(\cB_j)\ge j$.  The theorem would then give a quasi-greedy
Markushevich basis in a Banach space with no Schauder reordering.
\end{proof}

\section{Verification notes}

The proof has two load-bearing points.  First, a global greedy set really does
intersect each block in a block-greedy set; no comparison with coefficients in
other blocks is needed after this observation.  Second, at the global time
when the $m_j$th coordinate of a block appears, restriction of the global
prefix projection to $E_j$ is exactly the first-$m_j$ projection in the
induced block order.  Both are exact identities, not asymptotic estimates.

The uniform bound on the coordinate functionals prevents coefficient pathologies
across the direct sum and places the resulting system in the usual
semi-normalized, $M$-bounded framework.  No computational verification is
needed.  The separate file \texttt{verification.md} records an adversarial
check of the argument.

\section{Limitations and novelty status}

This is a reduction, not a construction of the finite-dimensional blocks.
The existence of a family satisfying the three hypotheses is open here.  No
converse is claimed: a dimension-free best-ordering bound for all finite
bases does not immediately produce a coherent order of an arbitrary infinite
basis.

For the bounded novelty check on 9 August 2026, the run registry and arXiv
were searched using combinations of the phrases \emph{finite-dimensional
quasi-greedy basis}, \emph{permutation}, \emph{best ordering constant},
\emph{Schauder reordering}, and \emph{direct sum}, together with arXiv ids
2510.06398 and 2510.13693.  No exact occurrence of the theorem above was
found.  The later paper \cite{AAB2025} proves a quasi-Banach counterexample
but explicitly states that the Banach restriction remains open.  Novelty of
this elementary reduction is plausible, not certified; it may be folklore.

\section{Human-review recommendation}

Review as a likely valid partial result.  The mathematical proof is short;
specialist effort is best spent on (i) checking whether the reduction is
already standard, and (ii) testing known finite-dimensional conditional
quasi-greedy constructions for divergence of $\beta$ under uniform $K$ and
$M$.

\begin{thebibliography}{9}
\bibitem{Berasategui2025}
M.~Berasategui,
\emph{Quasi-greedy Markushevich bases, duality and norming subspaces},
arXiv:2510.06398v1 (2025).

\bibitem{AAB2025}
F.~Albiac, J.~L.~Ansorena, and M.~Berasategui,
\emph{When Greedy Approximation Breaks: Counterexamples in Quasi-Banach
Spaces}, arXiv:2510.13693 (2025).
\end{thebibliography}

\end{document}
